\cotefile{157-4}
\scanpage{258}
\mspage{132}

que $Z(Y_0)$ est réduit à un point (les deux
facteurs $Z(Y_0)$, $Z(Y_1)$ étant $\ne\varnothing$). Cela
\deleted{\unsure{signifie}} \inserted{signifie} que $Y_0$ est un objet initial
dans $\mathrm{Hot}_{W}^{*}$. Il s'ensuit notamment
que $Y_0\amalg Y_0$ est aussi un objet initial,
et l'application can $Y_0\amalg Y_0\longrightarrow Y_0$ est donc
un iso dans $\mathrm{Hot}_{W}$, alors qu'\deleted{elle n'induit}
l'application induite sur les $\pi_0$ n'est pas
injective. (C'est $\pi_0(Y_0)\amalg\pi_0(Y_0)\unsure{-|\to}\pi_0(Y_0)$,
avec $\pi_0(Y_0)$ non vide.) Ce cas va être
traité ci-dessous.

Si dans \unsure{(a)|($*$)} $X$ est vide, \deleted{\unsure{les deux}}
$Y$ non vide (puisque $\pi_0(f)$ non \unsure{surj}),
alors $Y$ est objet initial en $\mathrm{Hot}_{W}$,
\unsure{et} et on voit \unsure{encore} que $Y\amalg Y_{\unsure{1|{}}}\unsure{-|\to}Y$
est iso dans $\mathrm{Hot}_{W}$, bien que \unsure{n'en donnant|n'en dérivant}
pas un iso \deleted{dans} des $\pi_0$. La prop. 9.\unsure{9|8}
va donc résulter du

\uline{Lemme} 9.\unsure{10}. Soit $W$ localiseur fondamental sans \deleted{\unsure{plus}}\inserted{\unsure{plus}} (axiomes $\mathrm{Loc}(1,\unsure{2|1},3)$).
Supposons qu'il existe $f\unsure{\in|:}\deleted{\unsure{X}}\!I\to\deleted{\unsure{Y}}\!J$ dans
\deleted{\unsure{$\mathrm{Hot}_{W}$}}\inserted{$\mathrm{Cat}$}, tel que \inserted{$\mathrm{Hot}_{W}(f)$ soit iso mais}
que $\pi_0(f)$ \uline{ne soit}
\uline{pas} \deleted{injectif} \inserted{\uline{surjectif}}. \deleted{\unsure{ainsi que}} Alors $\mathrm{Hot}_{W}^{\unsure{\wedge|*}}$
est équivalente à la catégorie ponctuelle \uline{$e$}, et
$\mathrm{Hot}_{W}$ est équivalente à $e$ (si $\varnothing\to e$ est un
iso dans $\mathrm{Hot}_{W}$) ou à $\Delta^1$ (dans le cas contraire)
% slanted
\margin{On \unsure{supp.|suppose} \unsure{tel} que $I$ soit $0$-connexe, et $J$ non $0$-connexe (i.e. $\mathrm{card}\,\pi_0(I)=1$, $\mathrm{card}\,\pi_0(J)\geq2$). \unsure{deuxième} \unsure{cas} \uuline{\unsure{traité}}}
% archive sheet number: 257
\scanpage{259}
\mspage{133}

La démonstration est celle donnée \deleted{pour} 9.7.
on prouve que si $\unsure{X|I},\unsure{Y|J}$ sont deux objets
de $\mathrm{Cat}$, \inserted{\uline{non vides}}, alors deux flèches
$f_0,f_1:X\mathrel{\unsure{=|\rightrightarrows}}Y$ sont égales dans $\mathrm{Hot}_{W}$,
par un argument d'homotopie, en
choisissant deux pts $i_0,i_1$ de $I$ qui
sont dans deux comp. connexes distinctes,
et tels que leurs images \unsure{$f_0,f_1$|$fi_0,fi_1$} dans $J$
soient dans la $\unsure{\overline{\mathrm m}|=}$ composante.

\uline{Prop. 9.11}. \deleted{La conjonction des} \inserted{Soit $W$ un localiseur fondamental}
(i.e. on a $\mathrm{Loc}(1,\unsure{2|1},3)$). Alors les conditions
suivantes sont équivalentes

a) On a $\mathrm{Loc}(5\deleted{\unsure{a}}\inserted{\mathrm{bis}})$ (axiome des \deleted{\unsure{diviseurs}}\inserted{pointes},
forme \deleted{\unsure{faible}}\inserted{forte}) et $\mathrm{Loc}(8)$ (axiome de connexité)

b) On a $\mathrm{Loc}.(5\,\mathrm{bis})$, et $\mathrm{Hot}_{W}$ n'est pas
équivalente à $e$ ou à $\Delta^1$ (cf. commentaire
9.8.)

c) Pour tt $f:X\to Y$, on a $f\in W$ sss
$\pi_0(f)$ iso, et \unsure{pour|après} \unsure{toute|tte} comp. \unsure{connexe|connexes}
$\unsure{X_i|X'}$ de $X$, si $Y_i$ désigne la comp correspondante de $Y$, \unsure{la flèche induite|les flèches induites} $f_i:X_i\to Y_i$
est dans $W$.

\deleted{De plus, ces conditions impliquent}

La raison \deleted{principale} pour laquelle je n'érige
pas cette propriété en un des \unsure{-} axiome
$\mathrm{Loc}(\unsure{\ldots|---})$, c'est que ce n'est \uline{pas} un axiome
de stabilité, mais un axiome de nature
\unsure{mixte}.
% archive sheet number: 258
\scanpage{260}
\mspage{134}

10. \uline{Application au dérivateur} $\mathrm{HOT}_{W}$

\uline{Théorème} 10.1. Soit $W$ un localiseur de $\mathrm{Cat}$ satisfaisant les conditions $\mathrm{Loc}(1\deleted{\unsure{\mathrm{bis}}},2,3,4,7\,\mathrm{bis})$
(axiome de \unsure{saturation|structure} forme forte\inserted{(\unsure{éventuellement})}, de l'objet unité,
de localisation forme faible, du carré cocartésien, de la
\deleted{des limites inductives et} $\unsure{L|\mathrm I}$-accessibilité). \deleted{\unsure{Alors}}
% slanted
\margin{\uline{NB} $\mathrm{Loc}(1\,\mathrm{bis})$ i.e. $W=W^{\unsure{s|g}}$, est déjà impliqué par la condition que $W$ stable \unsure{sous|par} \uline{lim} filtrantes}
Soit $\mathrm{HOT}_{W}$ le prédérivateur \inserted{sur $(\mathrm{Cat})$} associé à $(\mathrm{Cat},W)$
de la façon habituelle,

\[
 I\longrightarrow \deleted{\mathrm{Hot}}\unsure{\mathrm{Cat}|\mathrm{Hot}}(I\unsure{]|)}W(I)^{-1}
 =\underline{\mathrm{Hom}}(I,\mathrm{Cat}]W(I)^{-1}.
 \tag*{(10.1.1.)}
\]

On a alors ce qui suit.

\uline{1°}) \deleted{Les foncteurs $u_!$ existent} Pour toute flèche
$u:I\unsure{-|\to}J$ de $\mathrm{Cat}$, le foncteur $u_!$ (adjoint
à gauche de $u^*$) existe. (NB. on n'a pas
besoin pour ceci de $\mathrm{Loc}(4)$ ni de $\mathrm{Loc}(7\,\mathrm{bis})$,
il suffit que $W$ soit un localiseur fondamental),
et il se calcule de la façon standart
(axiomes $\mathrm{Der}\,4\,\mathrm{bis}$ et $\mathrm{Der}\,5\,\mathrm{bis}$). Le foncteur
$u_*$, adjoint à droite de $u^*$, existe \textsuperscript{*}[si
% slanted
\margin{\uuline{\unsure{faux|gros}} ! Cf $\overline{\underline{\mathrm{XIX}}}$ p. \unsure{28|29}}
$u$ est \deleted{une projection} de la forme $v\times\mathrm{id}_{S}$, où
$v:I_0\unsure{-|\to}J_0$ est une flèche entre catégories
ordonnées finies, et $S$ un objet quelconque
de $\mathrm{Cat}$,] \inserted{(et ce foncteur $u_*$ se calcule de façon standart (axiomes $\mathrm{Der}\,4$ et $\mathrm{Der}\,5$).)}
\unsure{D'où on a|Donc on a} une théorie des
carrés $W$-\unsure{cocartésiens|cartésiens} et des carrés $W$-cartésiens
dans les catégories $\mathrm{HOT}_{W}(S)$, et des suites
% archive sheet number: 259
\scanpage{261}
\mspage{133 bis}

\uline{Prop. 9.12} : Soit $W$ un localiseur fondamental (sans plus),
i.e satisfaisant $\mathrm{Loc}(1,2,3)$. \deleted{Pour que} Les conditions
suivantes sont équivalentes

a) $W$ satisfait l'axiome de connexité $\mathrm{Loc}(8)$, i.e.
$f\in W\Longrightarrow\pi_0(f)$ bijectif.

b) Si $f,g:X\mathrel{\unsure{=|\rightrightarrows}}Y$ dans $\mathrm{Cat}$ sont telles que $\mathrm{hot}_{W}(f)=\mathrm{hot}_{W}(g)$
alors $\pi_0(f)=\pi_0(g)$.

c) \deleted{Il \unsure{y a}} existe $X,Y$ dans $\mathrm{Cat}$ tel que \inserted{$\mathrm{hot}_{W}^{X,Y}:$} $\mathrm{Hom}_{\mathrm{Cat}}(X,Y)\to\mathrm{Hom}_{W}(X,Y)$ non \unsure{constante|constant}.

c') La sous-catégorie $\mathrm{Hot}_{W}^{\unsure{*|\wedge}}$ \deleted{\unsure{de}} de $\mathrm{Hot}_{W}$, formée des
objets \deleted{qui} non vides, n'est pas équivalente à la catégorie
ponctuelle.

c'') La catégorie $\mathrm{Hot}_{W}$ n'est pas équivalente à la
catégorie $e$ ni à la catégorie $\Delta^1$.

d\deleted{\unsure{'}}) L'\unsure{ens} des classes d'isomorphie d'objets de
$\mathrm{Hot}_{W}$ est infini.

\uline{Dém} On a $a\Longrightarrow b$, puisque a) équivaut à dire
que le foncteur $\pi_0:\mathrm{Cat}\to\mathrm{Ens}$ se factorise
par $\mathrm{Hot}_{W}$. D'ailleurs le foncteur correspondant

\[
 \mathrm{Hot}_{W}\longrightarrow(\mathrm{Ens})
\]

est ess. surjectif -- il y a un inverse à isom
près, $I\unsure{-|\to}\mathrm{hot}_{W}(I)$. Cela prouve que
a) implique \unsure{d')|d)}, et on a $\unsure{d'|d}\Rightarrow\unsure{c'|d'},\,\unsure{d'|d}\Rightarrow\unsure{c''|d''}$
tautologiquement. Un argument $\pm$ tautologique
montre d'ailleurs que si $\mathrm{Hot}_{W}^{*}$ est équiv. à $e$,
alors $\mathrm{Hot}_{W}$ est équivalente à $e$ ou à $\Delta^1$
(car $\varnothing$ est objet initial de $\mathrm{Hot}_{W}$) -- donc $c''\Rightarrow c'$.
Donc on a $\pm$ tautologiquement

\diagram{\begin{tikzcd}[ampersand replacement=\&,column sep=small,row sep=small]
a \arrow[r,Rightarrow] \arrow[d,Rightarrow] \& d \arrow[r,Rightarrow] \& c'' \arrow[r,Rightarrow] \& c' \\
b \arrow[rr,Rightarrow] \arrow[rrru,Rightarrow] \arrow[rrru,draw=black,decorate,decoration={snake,amplitude=1pt,segment length=7pt},no head,"\deleted{\vphantom{\Rightarrow}}"{inner sep=0pt}] \& \& c \arrow[ur,Rightarrow] \&
\end{tikzcd}}

où les implications $b\Rightarrow c\nearrow c'$ \deleted{et celle} \inserted{sont elles}
aussi triviales. \unsure{Donc|D'où}
% archive sheet number: 260
\scanpage{262}
\mspage{133 ter}

tout revient à prouver que $c)\Rightarrow a)$, ou encore :

\uline{Corollaire 9.13}. Si l'axiome de connexité n'est pas
satisfait, alors \deleted{deux flè} si $X,Y$ sont deux objets de $\mathrm{Cat}$,
\deleted{deux} \inserted{toutes les} flèches $f\deleted{\unsure{,g}}:X\mathrel{\unsure{\rightrightarrows|\to}}Y$ \inserted{de $\mathrm{Cat}$} sont \deleted{toujours} égales dans
$\mathrm{Hot}_{W}$ i.e. l'application $\mathrm{Hom}_{\mathrm{Cat}}(X,Y)\unsure{-|\to}\mathrm{Hom}_{W}(X,Y)$ est
constante.

Cela implique donc que \unsure{tous} les objets de $\mathrm{Cat}^*$ i.e.
tous les objets non vides de $\mathrm{Cat}^{\unsure{*|{}}}$ sont isomorphes
dans $\mathrm{Hot}_{W}$, et on trouve un isomorphisme can.
entre elles, \deleted{\unsure{D'où est de}} par transitivité. \unsure{Soit}
\deleted{\unsure{donne de la}} \inserted{alors $\mathscr H$} la catégorie \inserted{quotient de $\mathrm{Cat}$} ayant $\unsure{\overline{\mathrm m}|=}$ objets que $\mathrm{Cat}$
et ayant exactement une flèche entre deux objets
\unsure{non|deux} vides de $\mathrm{Cat}$, et exactement une flèche de $\varnothing$
dans tt autre objet, et pas de flèche d'un objet
non vide de $\mathrm{Cat}$ dans $\varnothing$). C'est donc une cat. équivalente
à $\Delta^1$. On a \deleted{\unsure{par construction une factorisation}} \inserted{vu que $\mathrm{Cat}\unsure{-|\to}\mathrm{Hot}_{W}$ se factorise}
en

\[
 \mathrm{Cat}\xrightarrow[\mathrm{can}]{}\mathscr H\longrightarrow\mathrm{Hot}_{W},
\]

\deleted{\unsure{qui est}} Si \deleted{\unsure{$\mathscr H$}} $\mathrm{Cat}\unsure{-|\to}\mathscr H$ transforme toute flèche \unsure{de $W$}
en iso, alors il est immédiat que $\mathscr H$ \deleted{satisfait} \inserted{est}
la \deleted{\unsure{solution}}\inserted{solution} \deleted{\unsure{au}}\inserted{au} problème universel définissant
$\mathrm{Hot}_{W}$, donc que $\deleted{\mathrm{Cat}\to}\mathscr H\longrightarrow\mathrm{Hot}_{W}$ est un iso.
Cela signifie aussi que \deleted{\unsure{$\mathscr H$}} $W$ ne contient pas
de flèches $\varnothing\to X$ avec $\unsure{\alpha|X}$ non vide (auquel
cas $W$ contient toutes \unsure{ces|les} flèches). Donc
dans ce cas $\mathrm{Hot}_{W}\approx\Delta^1$. Dans le cas \deleted{\unsure{inverse}}\inserted{contraire},
\deleted{\unsure{S}} soit $\mathscr H'$ la catégorie déduite de $\mathscr H$ en inversant
les flèches $\varnothing\longrightarrow X$, i.e. \unsure{le|la} groupoïde
1-connexe sur l'ens d'objets $\mathrm{Ob}\,\mathrm{Cat}$. Ce n'est
plus une catégorie quotient de $\mathrm{Cat}$, mais
on a encore une factorisation

\[
 \mathrm{Cat}\xrightarrow[\text{épi}]{}\mathscr H
 \mathrel{\mathop{\unsure{-|\to}}\limits_{\mathrm{loc.}}}\mathscr H'
 \longleftarrow\mathrm{Hot}_{W},
\]
% archive sheet number: 261
\scanpage{263}
\mspage{133 quater}

et cette fois c'est $\mathscr H'$ qui est solution du pb.
universel ($\mathrm{Cat}\unsure{-|\to}\mathscr H'$ transformant flèches de
$W$ en iso). Donc dans ce cas $\mathscr H'\unsure{-|\to}\mathrm{Hot}_{W}$ est
un iso, et $\mathrm{Hot}_{W}$ est équiv. à $e$. Dans les
deux cas, $\mathrm{Hot}_{W}^{*}$ est équivalente à $e$.

Donc, la prop. 9.12. résulte du cor. \unsure{9,13.|9.13.}, qu'il
s'agit à présent de démontrer.

Soit $f:\deleted{\unsure{X}}I\to\deleted{\unsure{Y}}J$ \inserted{de $\mathrm{Cat}$} telle que $\pi_0(f)$ ne soit pas injective.
On a vu dans la dém. de 9.7 comment on en
déduit (en choisissant deux éléments $i_0,i_1$ dans deux
comp. connexes distinctes de $I$, telles que
\deleted{\unsure{$j_1$}} $j_0=f(i_0)$ et $j_1=f(i_1)$ \unsure{'} soient dans une même
composante connexe de $\unsure{J|D}$) comment on déduit
que pour $X,Y$ dans $\mathrm{Cat}$, deux flèches $f_0,f_1:X\mathrel{\unsure{=|\rightrightarrows}}Y$
sont toujours égales dans \deleted{\unsure{$W$}} $\mathrm{Hot}_{W}$. (cf lemme 9.10.)

Soit maintenant

\[
 f:I\unsure{-|\to}J
\]

dans $\mathrm{Cat}$ telle que $\pi_0(f)$ ne soit pas surjective,
donc $J=J_0\amalg J_{\unsure{1|0}}$, avec $f$ se factorisant par
$J_0$, et $J_1\ne\varnothing$. Considérons alors la cat. discrète
à deux objets,

\[
 \unsure{\varepsilon|E}=\{0,1\},
\]

et deux \deleted{\unsure{appl}} flèches

\[
 g_0,g_1:J\rightrightarrows\varepsilon,
\]

$g_0$ étant \unsure{constante} de valeur $0$, et $g_1$ étant
égale à $0$ sur $J_0$, à $1$ sur $J_1$. On a

\[
 g_0f=g_1f\quad\text{\unsure{et|d}}
\]

donc, comme $f\in W$,

\[
 g_0=g_1\qquad\text{dans }\mathrm{Hot}_{W}.
\]

On en conclut, en prenant un point $j_1$ dans $J_1$,
et l'inclusion $\alpha_{\unsure{j_1|1}}:e\longrightarrow J$ correspondante, que
% archive sheet number: 262
\scanpage{264}
\mspage{133 quinquies}

\[
 g_0\alpha_{j_1}=g_1\alpha_{j_1}
\]

i.e. les deux sections de $\unsure{\varepsilon|E}$ sur $e$ sont égales
dans $\mathrm{Hot}_{W}$. Comme le foncteur $\mathrm{Cat}\unsure{-|\to}\mathrm{Hot}_{W}$
commute aux produits finis, il s'ensuit
que \inserted{pour $X$ dans $\mathrm{Cat}$,} les deux flèches

\diagram{\begin{tikzcd}[ampersand replacement=\&,column sep=large]
X \arrow[r,shift left=0.65ex,"{\unsure{\partial|\delta}}_0^X"] \arrow[r,shift right=0.65ex,"\partial_1^X"'] \& X\times\varepsilon=X\amalg X
\end{tikzcd}}

sont égales dans $\mathrm{Hot}_{W}$. Si alors

\[
 f_0,f_1:X\mathrel{\unsure{=|\rightrightarrows}}Y
\]

sont deux flèches dans $\mathrm{Cat}$, elles sont de
la forme

\[
 f_0=h\partial_0^X,\qquad f_1=h\partial_1^X,
\]

avec

\[
 h=(f_0,f_1):X\amalg X\longrightarrow\unsure{X|Y},
\]

et comme $\partial_0^{\unsure{\alpha|X}}=\partial_1^X$ dans $\mathrm{Hot}_{W}$, on en
conclut $f_0=f_1$ dans $\mathrm{Hot}_{W}$, qed.
% archive sheet number: 263
\scanpage{265}
\mspage{135}

longues de \deleted{\unsure{co}}suspension et de cosuspension dans
les catégories $\mathrm{HOT}_{W}^{\bullet}(S)$. (Plus exactement, ce sont
des théories relatives aux dérivateurs induits
$\mathrm{HOT}_{W,S}$, $\deleted{\unsure{\mathrm H}}\mathrm{HOT}_{W,S}^{\bullet}$ sur $S$, $I\unsure{-|\mapsto}\mathrm{HOT}_{W}(S\times I)$ resp.
$I\unsure{\mapsto|-}\mathrm{HOT}_{W}^{\bullet}(S\times I\ldots$) De plus, les carrés
cocartésiens et les carrés cartésiens relatifs
aux dérivateurs induits $\mathrm{HOT}_{W,S}$ sont
\uline{exacts}, i.e les suites \deleted{exactes} longues
de suspension et de cosuspension sont
exactes (axiomes d'exactitude à \deleted{gauche} \inserted{droite}
et à gauche $\mathrm{Der}\,6\,\mathrm{bis}$ et $\mathrm{Der}\,6$.)

2°) Le \deleted{pré}dérivateur $\mathrm{HOT}_{W}$ satisfait aussi
l'axiome $\mathrm{Der}\,1$ de \deleted{\unsure{somme directe}}\inserted{décomposition (finie $*$)}
% slanted
\margin{sans aucun axiome sur $W$ !}
(pour $I=\coprod_{\alpha}I_{\alpha}$),
\deleted{\unsure{Ainsi}} et l'axiome $\mathrm{Der}\,2$ de localisation.
De plus, $W$ est fortement saturé i.e. $f\in W\Longleftrightarrow\mathrm{hot}_{W}(f)$ iso.

3°) Concernant l'axiome $\mathrm{Der}\,3$ et ses variantes
(cf I, \unsure{pages|p.ges} 111--112, 125), on a \deleted{le résultat} ceci :

\deleted{\unsure{a)}} a) \deleted{La \unsure{signification} D-équivalences $W$}
% slanted
\margin{On n'a \deleted{pas} besoin ici \unsure{que} de $\mathrm{Loc}(1\,\mathrm{ter},\unsure{2},3)$ ($\mathrm{Loc}(1\,\mathrm{ter})$ est la saturation forte i.e. $\deleted{\unsure{a}}f\in W\Longleftrightarrow\mathrm{hot}_{W}(f)$ iso.}
Soit $f$ \inserted{$:S'\to S$} une flèche \unsure{de|dans} $\mathrm{Cat}$. Pour\textsuperscript{*}
\unsure{que} ce soit une $\mathrm{HOT}_{W}$-équivalence, i.e. pour
que le foncteur $f^{\unsure{*|!}}:\mathrm{HOT}_{W}(S)\to\mathrm{HOT}_{W}(S')$ induise une
\uline{équivalence} entre les sous-catégories pleines
formées des \deleted{\unsure{$f$}} objets \uline{constants} sur $S,S'$, il f. et
s. que $f$ soit dans $W$. \deleted{\unsure{On a}} L'énoncé correspondant dans les dérivateurs induits $\mathrm{HOT}_{W}(I)$
est \deleted{valable} également valable, pour $I\ne\varnothing$. Donc
à fortiori, $\mathrm{HOT}_{W}$ et $\mathrm{HOT}_{W}$ satisfont
% archive sheet number: 264
\scanpage{266}
\mspage{136}

l'axiome de l'objet \deleted{\unsure{unité|Cité}}\inserted{final}, i.e. \deleted{la $\mathrm{HOT}$} si
$S$ dans $\mathrm{Cat}$ a un objet final, alors $S\unsure{-|\to}e$
est une $\mathrm{HOT}_{W}$-équivalence, et \inserted{même} une
$\mathrm{HOT}_{W,I}$ équivalence pour tt $I$ dans $\mathrm{Cat}$.

b) Supposons que $W$ satisfasse
$\mathrm{Loc}(1\,\mathrm{ter},2,3,5\,\mathrm{b})$ (\unsure{donc|i.e.} l'axiome \inserted{\unsure{\mbox{Der 5 b}|restreint}} des $W$-fibrations, variante \unsure{1|b|2}), \deleted{\unsure{il}} \inserted{Alors pour toute}
$f\in W$, $\mathrm{HOT}_{W}^{\mathrm{lc}}(S)\unsure{-|\to}\mathrm{HOT}_{W}^{\mathrm{lc}}(S')$ est une
\uline{équivalence de catégories} (la réciproque
étant claire par a). Idem pour
$\mathrm{HOT}_{W,I}^{\mathrm{lc}}(\unsure{J|S})\unsure{-|\to}\mathrm{HOT}_{W,I}^{\mathrm{lc}}(S')$ (i.e. pour
$\mathrm{HOT}_{W}^{\mathrm{lc}/I}(S\times I)\longrightarrow\mathrm{HOT}_{W}^{\mathrm{lc}/\unsure{I'|I}}(S'\times I)$). En
d'autres termes, pour \inserted{\unsure{le dérivateur|les dérivateurs}} $\mathrm{HOT}_{W}$ et
% slanted
\margin{À vérifier avec soin dans le cas $\unsure{\mathbb D|D}=\mathrm{HOT}_{W,I}$}
pour tous les $\mathrm{HOT}_{W,I}$, \deleted{toute $\mathbb D$-}\inserted{($\mathbb D$ étant un de}
ces dérivateurs), toute $\mathbb D$-équivalence
est une $\mathbb D$-équivalence $2$-complète
i.e. pour $f:S'\unsure{-|\to}S$, pour que $f^*:\mathbb D(S)\to\mathbb D(S')$
induise une \deleted{\unsure{$\mathrm{HOT}$}} équivalence entre les
sous-catégories $\mathbb D^{\mathrm{lc}}(S)\unsure{-|\to}\mathbb D^{\mathrm{lc}}(S')$, il suffit \inserted{déjà}
qu'elle induise un foncteur pl. fid.
sur les sous-catégories pleines formées des
objets constants, et ceci équivaut (si $I\ne\varnothing$)
à la condition que $f\in W$.
% archive sheet number: 265
\scanpage{267}
\mspage{137}

\uline{Commentaire} : Dans cet énoncé récapitulatif de ce
que je sais prouver à présent pour $\mathrm{HOT}_{W}$, j'ai
un peu mélangé les provenances diverses --
ce qui est mis ``pour mémoire'' de $\overline{\underline{\mathrm{VI}}}$
(et qui ne dépend pas de $\mathrm{Loc}(4,7\ \text{ou }7\,\mathrm{bis})$),
i.e ce qui dépend pour l'essentiel de la
construction de $u_!$ donnée dans $\overline{\underline{\unsure{\mathrm{VII}|\mathrm{VI}}}}$, \deleted{et}
(donc ce qui concerne $u_!$ dans 1°, \uline{\unsure{moins|mais}}
l'exactitude à droite, et \unsure{tout} ce qui concerne
3° a) b)).

La démonstration de 10.1. résultera des
développements des § suivants, \deleted{ainsi que} \inserted{compte tenu} du
résultat principal du § 8 (Th 8.4.), et \unsure{de|du}
celui du § 7 (cor. 7.4\unsure{\textsuperscript{\ensuremath{\prime}}|.}), qui sont les
ingrédients ``durs'' du th. 10.1.
% archive sheet number: 266
\scanpage{268}

\margin{\unsure{3|$\rightharpoonup$}}
$\overline{\underline{\mathrm{XVI}}}$ \unsure{Localiseurs}
\unsure{fondamentaux}
$W\subset\mathrm{Fl}(\unsure{\mathrm{Cat}})$

\diagram{\(\begin{array}{ll@{\quad}l}
\mathrm{Loc}\ 1,1\,\mathrm{bis},1\,\mathrm{ter}&\text{Saturation}&
 \left.\vphantom{\begin{array}{c}1\\2\end{array}}\right\}\ \text{\unsure{§ 1}}\quad\text{p. }\unsure{1} \text{--} \unsure{8|7}\\
\mathrm{Loc}\ 2&\text{Objet final}&\\
\mathrm{Loc}\ 3,3\,\mathrm{bis}&\text{Localisation\unsure{;|,} !}&\\
\mathrm{Loc}\ 4\ (a,b,\unsure{M_0|4_0|40|4o})&\text{Sommes directes}&\text{§ 2}\ \}\quad\text{p. }8\text{--}15\\
\mathrm{Loc}\ 5&\text{Connexité}&\\
\mathrm{Loc}\ 6\ (a,b,c)&W\text{-cofibrations}&\left\{\begin{array}{l}\text{§ 3}\quad\text{p. }16\text{--}\unsure{27}\\[-2pt]\text{§ 4}\quad\text{p. }27\text{--}\unsure{70}\end{array}\right.\\
\mathrm{Loc}\ 7\ (a,b,c,\unsure{F|f})&W\text{-fibrations}&\text{§ 5}\ \}\quad71\text{--}\unsure{83}\\
\mathrm{Loc}\ 8,8\,\mathrm{bis}&\text{Limites}&\\
\mathrm{Loc}\ 9&\text{Accessibilité}&\\
[\,\mathrm{Loc}\ 10&\text{Descente (?).}&
\end{array}\)}

Axiome très fort, que
je n'ai pas explicité. Je sais
que \deleted{\unsure{vers}}\inserted{$\mathrm{Loc}(5,7\,\mathrm D_{W},10)$} \deleted{\illegible{3}} \inserted{\uline{caractérise} déjà}
$\unsure{\mathcal W_{\infty}|W_{\infty}}$, cf $\overline{\underline{\mathrm{XII}}}$ prouvant $\mathrm{Loc}\,7\,\mathrm D_{W}+\deleted{\unsure{\alpha}}$
$\Longrightarrow(W\subset\mathcal W_{\infty}\ \text{ou}\ W=\mathrm{Fl}(\mathrm{Cat}))$

\diagram{\(\begin{array}{ll@{\qquad}l}
1.&\mathrm{Loc}(1,2,3)&\text{p. }1\text{--}7\\
2.&\mathrm{Loc}(4,5)&\text{p. }8\text{--}\unsure{16|15}\\
\begin{array}{l}3.\\4.\end{array}&\left.\begin{array}{l}\mathrm{Loc}(6)\\\mathrm{Loc}(7)\end{array}\right\}&
\begin{array}{l}16\text{--}27\\27\text{--}70\end{array}\\
5.&\mathrm{Loc}(8,9)&71\text{--}\unsure{83}\\
6.&\mathrm{Loc}(\unsure{5|D},7\,\mathrm D_{W_W})\Rightarrow W\subset\mathcal W_{\infty}&\\
&\text{Commentaires (types d'homotopie}&83\text{--}96\\
&\text{simples etc)}&\\
\deleted{\unsure{7}}&\deleted{\text{Dépendance des axiomes}}&\\
7.&\text{Relations avec les axiomes sur}&97\text{--}123\\
&\text{les dérivateurs}&\\
8.&\text{Dépendances des axiomes}&
\end{array}\)}
% archive sheet number: 267
\scanpage{269}
\mspage{1}

Les localiseurs \unsure{fondamentaux}

1. \uline{Axiomes} $\mathrm{Loc}(1)$ \unsure{à|:} $\mathrm{Loc}(3)$ (et \unsure{variantes}) : \uline{localiseurs fondamentaux},

J'ai envie de reprendre l'exposé des axiomes
utiles pour un ``localiseur \unsure{fondamental}

\[
W\subset\mathrm{Fl}(\mathrm{Cat}),
\]

en me guidant cette fois \unsure{plus} sur la \unsure{signification}
géométrique des axiomes, \unsure{que} sur les
aléas $\pm$ techniques \unsure{qu'on|que je} rencontre \unsure{(ou ne rencontre pas)} pour les \unsure{certifier|vérifier}. \unsure{Je suivrai|Je suis} \unsure{donc}
dans cet esprit la présentation \unsure{de|du} $\unsure{\mathrm{XV}|\mathrm{X}}$ § 9.

\uline{Loc} 1 (\uline{saturation}) \quad a) $\mathrm{Iso}\subset W$

b) Si deux parmi $\unsure{f,g,gf|f,g,fg}$ dans $W$, aussi le \unsure{$3^{\mathrm{e}}$|$3^o$}

c) Si $gf$ et $fg$ \unsure{sont} dans $W$, $f$ et $g$ \unsure{aussi}.

\margin{\unsure{1|!}}Les conditions a) \inserted{et} b) tout \unsure{seuls|seul}, je vais les appeler
(pour un $M\unsure{\in|\subset}\mathrm{Fl}\,M$ général) \uline{saturation faible} du
localiseur. Les conditions a) b) c), \uline{saturation modérée}
(``mild saturation''). \unsure{Autres variantes|Autre variante}

\uline{Loc} $(1\,\mathrm{bis})$. Les conditions a) b) c) de Loc 1, plus

d) $W$ stable par facteurs directs.

C'est bien sûr impliqué par :

\uline{Loc} $(1\,\mathrm{ter})$ $W$ est \uline{fortement saturé}, i.e. si $\unsure{f\in\mathrm{Fl}\,\mathrm{Cat}|f\in\mathrm{Cat}}$ est
telle que $\mathrm{hot}_{W}(f)$ \unsure{iso}, alors $f\in W$.

\deleted{\unsure{2|L}} \uline{Loc} $\unsure{(2)}$ (\uline{Objet final}). Si $X$ dans $\mathrm{Cat}$ a un \unsure{obj.|objet} final,
alors $X$ est $W$-asphérique i.e. $X\longrightarrow e$ est dans $W$.

Parfois, cet axiome est conséquence de la variante
affaiblie suivante

\uline{Loc} $\unsure{(2_0)|(2_o)}$ (\unsure{Axiome} du segment) $\Delta^1$ est $W$-asphérique
% archive sheet number: 268
\scanpage{270}
\mspage{2}

\uline{Loc} 3 (\uline{localisation}, forme faible) : Soit $f:X\longrightarrow S$ une
flèche \unsure{dans|de} $\mathrm{Cat}$ $W$-asphérique, i.e. telle que
$\forall\,s\in\mathrm{Ob}\,S$, \inserted{la cat. localisée $X_{/s}$ de $X$ \unsure{sur|en} $s$} $X_{/s}=X\times_S S_{/s}$ soit $W$-asphérique. Alors $f\in W$.

% slanted
\margin{Dans une version finale, il faudrait intervertir les notations pour $\mathrm{Loc}(3)$ et pour $\mathrm{Loc}(3\,\mathrm{bis})$. Car ce dernier m'apparaît comme le plus important des deux -- il est \unsure{obligé|satisfait}, pour \unsure{obtenir|tous} les localiseurs \deleted{\unsure{qui}} \unsure{ayant|peuvent} \unsure{été|être} utilisés. \unsure{Ceci|Ce} au moins (\unsure{sauf erreur}), est \deleted{\unsure{va}}valable pour tt localiseur de la forme $W_{\unsure{\mathbb D|D}}$, où $\mathbb D$ est un dérivateur (ou seulement un prédérivateur satisfaisant $\mathbb D\unsure{\mathrm{er}}\,(4,5)$).}

\uline{Rappels}. 1. $\mathrm{Loc}(1,2,3)$ impliquent que l'on a

\[
W=W^o\qquad\text{i.e}\quad f\in W\iff f^o\in W,
\]

\unsure{dénotant|dénoté} par $X\longmapsto X^o$, $f\longmapsto f^o$ le
foncteur de passage à la catégorie opposée, \inserted{\uline{foncteur}} qui

\[
X\longmapsto X^o:\quad \mathrm{Cat}\xrightarrow{\sim}\mathrm{Cat}
\]

qui est un \uline{automorphisme} de $\mathrm{Cat}$ (donc
commute à toutes $\underrightarrow{\lim}$ et $\underrightarrow{\lim}$). Il
s'ensuit que l'ensemble d'axiomes
$\mathrm{Loc}(1,2,3)$ équivaut à $\mathrm{Loc}(1,2_{\unsure{0|o}}^{\mathrm{opp}},3^{\mathrm{opp}})$, où
$\mathrm{Loc}(2^{\mathrm{opp}})$ est l'axiome ``dual'' de $\mathrm{Loc}(2)$,
savoir que \deleted{tt \unsure{catégorie}} objet de $\mathrm{Cat}$ ayant
un objet initial est $W$-asphérique, et
\unsure{de m\textsuperscript{e}|de =}
$\mathrm{Loc}(3^{\mathrm{opp}})$ l'axiome que \deleted{\unsure{tout}} $f:X\to S$ \uline{$W$-coasphérique}, i.e telle que les \inserted{catégories colocalisées}
$ {}_{s\backslash}X$ \inserted{de $X$ sous $s$} \deleted{\unsure{sont}} \deleted{\unsure{soient}} $W$-asphériques, est dans $W$.

\deleted{\unsure{Loc (3 bis)}} \deleted{\unsure{mutuel}}
\uline{Définition} 1.1. On appelle \uline{localiseur fondamental}
une partie $W$ de $\mathrm{Fl}(\mathrm{Cat})$ satisfaisant les
axiomes $\mathrm{Loc}(1,2,3)$ (i.e. saturation modérée,
axiome de l'objet final, axiome de localisation).
% archive sheet number: 269
\scanpage{271}
\mspage{3}

\deleted{\uline{\unsure{Proposition|Définition}}}\inserted{Rappel} 1.2. \deleted{\unsure{sous la satisfaction de|sous la condition de}} Si $W$ est un
localiseur fondamental, alors tout \uline{homotopisme}
$f:X\to Y$ dans $\mathrm{Cat}$ est dans $W$ (c'est pour cela
essentiellement que l'on a imposé c) dans \inserted{la saturation modérée} $\mathrm{Loc}(1)$),
si $f,g:X\unsure{\rightrightarrows|\substack{\to\\-}}Y$ sont homotopes (i.e. appartiennent
à une \unsure{$=$} composante connexe de $\underline{\mathrm{Hom}}(X,Y)$),
alors $\mathrm{hot}_W(f)=\mathrm{hot}_W(g)$ et $f\in W\iff g\in W$.

\deleted{De plus} Voici la forme forte de $\mathrm{Loc}(3)$

\deleted{Proposition}

\uline{Loc} $(3\,\deleted{\unsure{\mathrm{ter}|\mathrm{bis}}}\inserted{\mathrm{bis}})$ (\uline{Localisation}, forme forte) : Soit $S$ dans $\mathrm{Cat}$
et $f:X\to Y$ un morphisme dans $\mathrm{Cat}_{/S}$, tel que
``$f$\unsure{''|$ {}^{a}$} soit une $W$-équivalence localement sur $S$'', i.e.
pour tout $s\in\mathrm{Ob}\,S$, le morphisme induit

\[
f_{/s}:X_{/s}\longrightarrow Y_{/s}
\]

pour les catégories localisées de $X,Y$ sur $S$ en $s$,
soit dans $W$. Alors $f\in W$.

L'axiome $\mathrm{Loc}(3)$ n'est autre que le cas particulier de $\mathrm{Loc}(3\,\mathrm{bis})$ relatif au cas où $Y=S$.

\uline{Proposition} 1.3. Soit $W$ \deleted{satisfaisant $\mathrm{Loc}(1)$}\inserted{$\subset\mathrm{Fl}(\unsure{M|\mathcal M})$}. Alors
on a les équivalences suivantes

\[
\begin{aligned}
\mathrm{Loc}(1,2,3\deleted{\unsure{\mathrm{bis}}})&\iff\mathrm{Loc}(1,2_0,3'\deleted{\unsure{\mathrm{bis}}})\\
\mathrm{Loc}(1,2,3\,\mathrm{bis})&\unsure{\Longleftarrow|\iff}\mathrm{Loc}(1,2_0,3'\,\mathrm{bis})
\end{aligned}
\]

où \deleted{\unsure{le}} les axiomes $\mathrm{Loc}(3'\deleted{\unsure{\mathrm{bis}}})$ \inserted{et $\mathrm{Loc}(3'\,\mathrm{bis})$} se formulent ainsi :

\deleted{\uline{Loc} $(3'\,\mathrm{bis})$ (\uline{Axiome \inserted{\unsure{direct|disant}} des Cat-fibrations}, forme \deleted{\unsure{élaguée|forte}} \inserted{faible}) : Soit
$S$ dans $\mathrm{Cat}$, $X,Y$ deux catégories fibrées scindées
sur $S$, $f:X\to Y$ un $S$-foncteur cartésien
strict \unsure{(i.e. compatible avec les scindages)}. Supposons que pour tout $s\in\mathrm{Ob}\,S$, le foncteur}
% archive sheet number: 270
\scanpage{272}
\mspage{P 31}

\deleted{entre le formel et l'algébrique (avec, comme deuxième source d'inspiration,
l'article fondamental GAGA de Serre). Quant à la contribution de Serre à laquelle
je fais allusion dans le texte, il s'agit bien sûr, avant tout, de l'introduction
par lui, en géométrie algébrique abstraite\unsure{,} du point de vue des f\,a\,i\,s\,c\,e\,a\,u\,x
(introduit par J\,e\,a\,n L\,e\,r\,a\,y une douzaine d'années auparavant, dans un
contexte tout différent), dans cet autre article fondamental déjà cité FAC
(``Faisceaux algébriques cohérents'').}

\deleted{A la lumière de ces ``rappels'', si je devais nommer les ``ancêtres'' immédiats
de la nouvelle vision géométrique, ce sont les noms de O\,s\,c\,a\,r Z\,a\,r\,i\,s\,k\,i,
A\,n\,d\,r\,é W\,e\,i\,l, J\,e\,a\,n L\,e\,r\,a\,y et J\,e\,a\,n-P\,i\,e\,r\,r\,e S\,e\,r\,r\,e
qui s'imposent à moi aussitôt. Parmi eux Serre a joué un rôle à part, du fait
que c'est par son intermédiaire surtout que j'ai eu connaissance non seulement
de ses propres idées, mais aussi des idées de Zariski, de Weil et de Leray qui
ont eu à jouer un rôle dans l'éclosion et dans le développement de la géométrie
nouvelle.}

\rule{7em}{0.4pt}

\deleted{Leurs idées, bien sûr, ont joué un rôle important lors de mes premiers pas dans
l'édification de la géométrie arithmétique\quad ; plus, il est vrai, comme points de
départ et comme o\,u\,t\,i\,l\,s (qu'il m'a fallu refaçonner plus ou moins de
toutes pièces, pour les besoins d'un contexte beaucoup plus vaste), que comme
une source d'inspiration qui aurait continué à nourrir\handwritten{\unsure{$ {}^{\sim}$}} mes rêves et mes projets,
au cours des mois et des années. De toutes façons, il était bien clair d'emblée
que, même refaçonnés, ces outils étaient très en deça de ce qui était requis,
pour faire même les tout premiers pas en direction des fantastiques conjectures.}

\deleted{11.\quad Les deux idées-forces cruciales dans le démarrage et dans le développement
de la géométrie nouvelle, ont été celle de s\,c\,h\,é\,m\,a et celle de t\,o\,p\,o\,s.
Apparues à peu près simultanément et en étroite symbiose l'une avec l'autre (*),
elles ont été comme un seul et même n\,e\,r\,f m\,o\,t\,e\,u\,r dans l'essor spectaculaire de la nouvelle géométrie, et ceci dès l'année même de leur apparition.
Pour terminer ce tour d'horizon sur mon oeuvre, il me reste à dire quelque mots
au sujet tout au moins de ces deux idées-là.}

\deleted{La notion de schéma est la plus naturelle, la plus ``évidente'' imaginable,
pour englober en une notion unique la série infinie de notions de ``variété''}

\rule{7em}{0.4pt}

\deleted{(*) Il est question de ce démarrage, qui se place en 1958, dans la note de b. de
p. (*) page 23 . La notion de\unsure{,} site ou de ``t\,o\,p\,o\,l\,o\,g\,i\,e d\,e
G\,r\,o\,t\,h\,e\,n\,d\,i\,e\,c\,k'' (version provisoire de celle de topos) est apparue
dans le sillage immédiat de la notion de schéma. C'est elle à son tour qui fournit le langage nouveau de la ``localisation'' ou de ``la descente'', utilisé à chaque
pas dans le développement du thème et de l'outil schématiques. La notion plus
intrinsèque et plus géométrique de t\,o\,p\,o\,s, restée d'abord implicite au cours
des années suivantes, se dégage surtout à partir de 1963, avec le développement
de la cohomologie étale, et s'impose peu à peu à moi comme la notion la plus fondamentale.}
% archive sheet number: 271
